% year: תשע"ט
% semester: ב
% moed: ב
% number: 1ב
% points: 10
% categories: func-limits, func-analysis
% source: exams/מבחנים/תשעט/חודיסמן מועד ב תשעט.pdf

\begin{question}
(10 נק') מצאו את כל האסימפטוטות של הפונקציה: $f(x)=\dfrac{x\sqrt{x^2-1}}{2x^2-1}$.
\end{question}

\begin{hint}
שימו לב לתחום ההגדרה $|x|\ge1$ ולכך ש-$\sqrt{x^2}=|x|$.
\end{hint}

\begin{solution}
\textbf{תחום הגדרה:} צריך $x^2-1\ge0$, כלומר $|x|\ge1$, וגם $2x^2-1\ne0$, כלומר $x\ne\pm\frac1{\sqrt2}$ – אך נקודות אלה ממילא אינן בתחום. לכן התחום הוא $(-\infty,-1]\cup[1,\infty)$.

\textbf{אסימפטוטות אנכיות:} בתחום המכנה מקיים $2x^2-1\ge1>0$, ולכן $f$ רציפה בכל תחומה (גם בקצוות $x=\pm1$, שם $f(\pm1)=0$). אין נקודה שבה $f$ שואפת לאינסוף, ולכן אין אסימפטוטות אנכיות.

\textbf{אסימפטוטות אופקיות:} מכיוון ש-$\sqrt{x^2-1}=|x|\sqrt{1-\frac1{x^2}}$,
\[ f(x)=\frac{x|x|\sqrt{1-\frac1{x^2}}}{x^2\left(2-\frac1{x^2}\right)}=\frac{|x|}{x}\cdot\frac{\sqrt{1-\frac1{x^2}}}{2-\frac1{x^2}}. \]
כאשר $x\to+\infty$: $\frac{|x|}{x}=1$ ולכן $f(x)\to\frac12$. כאשר $x\to-\infty$: $\frac{|x|}{x}=-1$ ולכן $f(x)\to-\frac12$. מכיוון שהגבולות סופיים, אין אסימפטוטות משופעות נוספות.

\textbf{תשובה:} $y=\frac12$ (ב-$+\infty$) ו-$y=-\frac12$ (ב-$-\infty$); אין אסימפטוטות אנכיות.
\end{solution}
